% TOP_PROBLEM: {"problemId": "problem.artin-s-primitive-root-conjecture", "canonicalTitle": "Artin's primitive root conjecture", "releaseRank": 197, "primaryDomain": "number_theory_arithmetic_geometry", "primaryDomainLabel": "Number theory & arithmetic geometry", "displayStatus": "open", "releaseStatus": "open"}
% !TeX program = lualatex
% Definition and sources only; this file is not an LLM solution attempt.
\documentclass[11pt]{article}
\usepackage[margin=25mm]{geometry}
\usepackage{fontspec}
\setmainfont{Latin Modern Roman}
\usepackage{amsmath,amssymb,mathtools}
\usepackage{unicode-math}
\setmathfont{Latin Modern Math}
\usepackage{xurl,hyperref}
\hypersetup{colorlinks=true,urlcolor=blue}
\setcounter{secnumdepth}{0}
\setlength{\parindent}{0pt}
\setlength{\parskip}{5pt}
\setlength{\emergencystretch}{3em}
\begin{document}
\section*{\texorpdfstring{Artin's primitive root conjecture}{Artin's primitive root conjecture}}\label{197}
\subsection{Definitions and mathematical statement}
For a prime $p\nmid a$, let $\operatorname{ord}_p(a)$ be the least $m>0$ with $a^m\equiv1\pmod p$. Then $a$ is a primitive root modulo $p$ exactly when $\operatorname{ord}_p(a)=p-1$. For each integer $a\ne-1$ which is not a square, the target requires a positive prime-relative density
\[
 \lim_{x\to\infty}\frac{\#\{p\le x:\operatorname{ord}_p(a)=p-1\}}{\pi(x)}=\delta(a)>0.
\]
The Artin density is expressed in the usual inclusion-exclusion convention by
$\delta(a)=\sum_{m\ge1}\mu(m)/[{\mathbb{Q}}(\zeta_m,a^{1/m}):{\mathbb{Q}}]$, equivalently its corrected Artin-product form. The algebraic degrees encode dependencies between prime-power obstructions. The density is among primes, not among all positive integers.
\par\smallskip\textit{Formulation and concept references:} \href{https://www.sciencedirect.com/science/article/pii/S0022314X24000829}{[S1]}.

\subsection{Short English statement}
Does every admissible integer generate the multiplicative group modulo a positive proportion of primes, with precisely Artin's predicted density?
\subsection{Sources}
\begin{itemize}
\item[S1]Journal of Number Theory 262 (2024), 161-185.
\url{https://www.sciencedirect.com/science/article/pii/S0022314X24000829}.
\textit{Formulation source.}
\end{itemize}
\end{document}
